\section{Conexión con los episodios anteriores: Agent, Context y lenguaje}

Este episodio forma, junto con los anteriores, una línea principal muy clara en la capa de aplicación. EP139 trata la historia técnica de los Agent y enfatiza la llamada a herramientas, la disolución de fronteras y la irradiación social; EP135 trata Elys y los dobles digitales, y enfatiza la obtención y el flujo de context y la red personificada; EP133, desde la perspectiva de los modelos del mundo, recuerda que más allá del lenguaje existe el mundo físico. El lugar de EP130 consiste en devolver estos conceptos técnicos a la mesa de trabajo del gerente de producto.

\subsection{Una misma palabra, significados en distintos niveles}

\begin{longtable}{p{0.22\textwidth}p{0.34\textwidth}p{0.35\textwidth}}
\toprule
Concepto & Énfasis en los episodios anteriores & Significado de producto en EP130 \\
\midrule
Agent & Llamada a herramientas, ejecución de tareas, cambios en los límites del sistema & Entra por la expresión en PPT y la creación, y se convierte en un objeto de productividad interactivo. \\
Context & Información personal, redes sociales, flujo de datos de los dobles digitales & El gerente de producto debe diseñar el contexto visible para el modelo, no solo dibujar flujos. \\
Lenguaje & En EP133 se discute como insuficiente para abarcar el mundo físico & En EP130 se considera el núcleo de gobierno del Agent y la puerta de entrada de la intención del usuario. \\
Amigo de IA & EP135 se inclina más hacia las relaciones y las redes sociales & EP130 enfatiza ofrecer primero valor práctico y, poco a poco, generar la sensación de amistad. \\
\bottomrule
\end{longtable}

\begin{knowledgebox}{Por qué importa esta conexión}
Si se mira solo EP130, parece una metodología de producto; puesto dentro de la serie completa, responde a «cómo un Agent se convierte en un producto que la gente común usa a largo plazo». La historia técnica, el flujo de context, los límites del lenguaje y la percepción del producto son distintos niveles de un mismo problema.
\end{knowledgebox}

\subsection{Resumen de la sección}

EP130 lleva al Agent del relato técnico de vuelta al relato de producto: no solo pregunta si el modelo puede hacerlo, sino también por qué el usuario lo recuerda, por qué ya no puede volver atrás y por qué está dispuesto a dejarlo entrar en su trabajo y en su red de relaciones.
